% !PRINTED 307
\weberpage{307}

F\"ur die erste Annahme $(a, b, c) = (1, 0, 0)$ wird die Determinante
 von $\sigma$ gleich $b_1 c_2 - c_1 b_2$, was auf sechs verschiedene
Arten $= 1$ werden kann, n\"amlich:

\[
\begin{pmatrix} b_1, & c_1 \\ b_2, & c_2 \end{pmatrix} = \begin{pmatrix} 1, & 0 \\ 0, & 1 \end{pmatrix}, \begin{pmatrix} 1, & 0 \\ 1, & 1 \end{pmatrix}, \begin{pmatrix} 0, & 1 \\ 1, & 0 \end{pmatrix}, \begin{pmatrix} 0, & 1 \\ 1, & 1 \end{pmatrix}, \begin{pmatrix} 1, & 1 \\ 1, & 0 \end{pmatrix}, \begin{pmatrix} 1, & 1 \\ 0, & 1 \end{pmatrix}.
\]

Dann k\"onnen noch $a_1, a_2$ auf vier verschiedene Arten angenommen
 werden. Also hat man, wenn die erste Zeile $(1, 0, 0)$
ist, 24 verschiedene Formen von $\sigma$. Dieselbe Zahl ergiebt sich
f\"ur die Annahme der ersten Zeile in den Formen $(0, 1, 0), (0, 0, 1)$.
Das Gleiche erh\"alt man aber auch f\"ur die anderen Annahmen
\"uber die erste Zeile. Denn ist diese z.\,B. $(0, 1, 1)$, so muss

\[
a_1 (b_2 - c_2) - a_2 (b_1 - c_1) \equiv 1 \text{ (mod 2)}
\]

sein, was wieder sechs M\"oglichkeiten f\"ur $a_1, b_1 - c_1, a_2, b_2 - c_2$
ergiebt, und da man $c_1, c_2$ auf vier Arten annehmen kann, so
erh\"alt man wieder 24 M\"oglichkeiten. Endlich muss, wenn die
erste Zeile $(1, 1, 1)$ ist,

\[
(a_1 - c_1)(b_2 - c_2) - (a_2 - c_2)(b_1 - c_1) \equiv 1 \text{ (mod 2)}
\]

sein, was zu derselben Zahl f\"uhrt. Die Gesammtzahl aller verschiedenen
 Substitutionen $\sigma$, d.\,h. der Grad von $S$ ist also 168.
Die Zahl 168 als Gradzahl einer Gruppe ist uns schon einmal
begegnet, n\"amlich bei der Gruppe $L_7$ (\S.\,66, 72), und es liegt
daher die Vermuthung nahe, dass die Gruppen $S$ und $L_7$
isomorph sein m\"ochten. Wenn sich diese Vermuthung best\"atigt,
so w\"urde die Untersuchung der Gruppe $S$ dadurch wesentlich
 erleichtert sein, dass wir die Divisoren der
Gruppe $L_7$ schon kennen.

Diese Vermuthung zu pr\"ufen, dient aber das Theorem I.
in \S.\,72.

Wir suchen zu diesem Zweck erzeugende Elemente $\chi, \omega, \Theta, \tau$
der Gruppe $S$ so auszuw\"ahlen, dass sie den charakteristischen
Bedingungen des Theorems I., \S.\,72 gen\"ugen. Dies kann auf
mehrfache Weise geschehen. Es fehlt freilich an einem allgemeinen
 Verfahren dazu, gelingt aber leicht durch einige Versuche.

Man w\"ahlt zun\"achst ein Element 7\textsuperscript{ten} Grades beliebig aus
und bildet daraus die Periode, etwa
%CONT>>
